% ===================================================================
%  TAFEL Nr. 2 — Der menschliche Körper. --- Pause. --- Spiele.
%  Traduction 2026 d'apres texto/io, controlee sur texto/fr et
%  texto/en. Consigne : voir texto/de/10-tafel-01.tex.
%
%  CETTE PAGE DONNE A LA SERIE SA TROISIEME LANGUE PLANIFIEE, ET
%  C'EST LA PLUS ANCIENNE DES TROIS.
%
%  Le tableau 1 romanche avait montre une ORTHOGRAPHE decidee en
%  1982 ; le tableau 1 estonien, des MOTS decides vers 1920. Le
%  vocabulaire de la seconde moitie de cette page-ci a ete decide en
%  1811, par un seul homme, et il est encore en place.
%
%  FRIEDRICH LUDWIG JAHN A FABRIQUE LA LANGUE DU GYMNASE. Il a
%  refuse « Gymnastik », qui etait grec et passait par le francais,
%  et fait « Turnen » ; puis « Turnhalle », « Turnlehrer »,
%  « Turngerät », et le nom de chaque engin :
%     « Barren »   les barres paralleles         (11)
%     « Reck »     la barre fixe
%     « Pferd »    le cheval d'arcons            (7)
%     « Ringe »    les anneaux                   (3)
%     « Barlauf »  les barres, le jeu de prisonnier (38)
%  Ces mots ne viennent de nulle part : « Barren » et « Reck » sont
%  des mots dialectaux qu'il a detournes, « Barlauf » un compose
%  qu'il a fait. L'alinea 21 de cette page est donc un inventaire de
%  1811 grave sur une planche francaise de 1926, et l'allemand le rend
%  sans un mot d'ecart.
%
%  ET L'ON VOIT OU LA DECISION S'ARRETE. Jahn a gagne sur tout ce qui
%  etait installe avant lui, et rien du tout sur ce qui est arrive
%  apres : l'alinea 28 nomme le Tennis sans le traduire,
%  l'alinea 22 le Boxen, l'alinea 22 encore le
%  Florett et le Trapez. Le programme etait clos
%  quand ces jeux-la ont traverse la Manche. UNE LANGUE DECIDEE NE
%  RESTE DECIDEE QUE LE TEMPS QUE DURE CELUI QUI DECIDE.
%
%  UNE SEULE EXCEPTION, ET ELLE A UNE DATE ELLE AUSSI : le
%  Fußball de l'alinea 30 est un calque, non un emprunt, et
%  le Fahrrad de l'alinea 31 en est un autre — il a remplace
%  « Veloziped » dans les annees 1890, sous la pression de
%  l'Allgemeiner Deutscher Sprachverein. Ce n'est plus Jahn, c'est
%  une seconde campagne, quatre-vingts ans apres, et elle n'a tenu que
%  sur deux mots.
%
%  LE TABLEAU 2 ESTONIEN AVAIT TROUVE « jalgpall » ET « jalgratas »,
%  LE MEME BALLON ET LE MEME VELO, COMPOSES DE LA MEME FACON. Trois
%  langues sans rapport — l'allemand, l'estonien, le romanche avec
%  son « ballape » — tombent sur le meme compose pour la meme chose.
%  Ce n'est pas la planification qui l'explique : un ballon qu'on
%  frappe du pied SE DECRIT, et qui le decrit tombe sur le pied et
%  sur la balle. La planification n'explique que le refus d'emprunter,
%  jamais la forme du remplacant.
%
%  LA PREMIERE MOITIE, LE CORPS, NE DONNE AUCUN EMPRUNT : Kopf,
%  Schädel, Gehirn, Haar, Gesicht, Stirn, Auge, Nase, Mund, Kinn,
%  Zunge, Hals, Brust, Herz, Rippen, Rücken, Bauch, Magen, Arm, Bein,
%  Schulter, Ellbogen, Hand, Finger, Nagel, Ader, Blut, Knie, Wade,
%  Haut, Fuß, Zehen, Ferse. C'est exactement ce que l'estonien avait
%  donne, et pour la meme raison : on emprunte le nom d'un objet
%  qu'un marchand apporte, on n'emprunte pas le nom de sa propre main.
%  Deux mots seulement font exception, et ce sont les deux que l'ecole
%  a apportes : Muskel et Bizeps, a l'alinea 15.
% ===================================================================

%%K t02-apar-1 apar
\VUsaut{4.0mm}
\VUtitre{15.0pt}{60}{TAFEL Nr. 2}
\VUsaut{2.0mm}
\VUtitre{11.4pt}{0}{Der menschliche Körper. --- Pause.}
\VUtitre{11.4pt}{0}{Spiele.}

%%K t02-tit-0 sub
\VUtitre{13.2pt}{0}{Der menschliche Körper.}
\VUsaut{2.4mm}
\VUcentre{10.2pt}{0}{\textit{(Eine einfache Naturkundestunde.)}}
\VUsaut{3.0mm}

%%K t02-01-1 p
1. --- Die wichtigsten Teile des menschlichen Körpers sind: der \VUgras{Kopf}, der
\VUgras{Rumpf} und die \VUgras{Gliedmaßen}.

%%K t02-tit-1 sub
\VUpk{10.2pt}{40}{Der Kopf.}
\VUsaut{3.0mm}

%%K t02-02-1 p
2. --- Der Kopf besteht aus zwei Teilen: dem \VUgras{Schädel}, der das
\VUgras{Gehirn} umschließt und von den \VUgras{Haaren} bedeckt ist, und dem
\VUgras{Gesicht}.

%%K t02-03-1 p
3. --- Auf unserer Zeichnung sehen wir weder den Schädel noch das Gehirn; die
wichtigen Teile des Gesichts aber sehen wir sehr deutlich.

%%K t02-04-1 p
4. --- Zuerst kommt die \VUgras{Stirn}, die auf einen kräftigen Verstand und
einen stets tätigen Geist schließen lässt. Die \VUgras{Schläfen} sind sehr offen.
Das \VUgras{Auge} ist lebhaft und weit geöffnet. Eine dichte \VUgras{Augenbraue}
und lange \VUgras{Wimpern} schützen es.

%%K t02-05-1 p
5. --- Dann kommen die \VUgras{Nase} und die \VUgras{Nasenlöcher}. Mit der Nase
riechen und atmen wir. Zu beiden Seiten der Nase liegen die \VUgras{Wangen} und
weiter außen die \VUgras{Ohren}. Der untere Teil des Gesichts besteht aus dem
\VUgras{Mund} und dem \VUgras{Kinn}.

%%K t02-06-1 p
6. --- Sie sind kaum zu sehen, weil der \VUgras{Schnurrbart} und der \VUgras{Bart}
sie verdecken. Wir wissen aber, dass der Mund die beiden \VUgras{Lippen}, den
\VUgras{Oberkiefer} und den \VUgras{Unterkiefer} mit ihren \VUgras{Zähnen} und die
\VUgras{Zunge} enthält. Mit dem Mund sprechen und essen wir. Mit der Zunge
schmecken wir unsere Speisen.

%%K t02-07-1 p
7. --- Auge, Ohr, Nase und Mund sind, zusammen mit den Fingern, die Organe der
fünf Sinne: \VUgras{Sehen}, \VUgras{Hören}, \VUgras{Riechen}, \VUgras{Schmecken}
und \VUgras{Tasten}.

%%K t02-08-1 p
8. --- Der Kopf ist also einer der interessantesten Teile des menschlichen
Körpers.

%%K t02-tit-2 sub
\VUsaut{4.4mm}
\VUpk{10.2pt}{40}{Der Rumpf.}
\VUsaut{3.0mm}

%%K t02-09-1 p
9. --- Der \VUgras{Hals} verbindet den Kopf mit dem Rumpf. Zu ihm gehören die
\VUgras{Kehle} und der \VUgras{Nacken}.

%%K t02-10-1 p
10. --- Auch im Rumpf liegen viele lebenswichtige Organe. In der \VUgras{Brust}
sind die \VUgras{Lungen}, mit denen wir atmen, und das \VUgras{Herz}, das der
Mittelpunkt des Lebens ist. Wenn das Herz aufhört zu schlagen, stirbt das
Lebewesen.

%%K t02-11-1 p
11. --- Die \VUgras{Rippen} stützen die Brust.

%%K t02-12-1 p
12. --- Die übrigen Teile des Rumpfes sind die \VUgras{Seite}, der \VUgras{Rücken},
die \VUgras{Schultern} und der \VUgras{Unterleib} (der Bauch). Zum Schlafen legen
wir uns auf die Seite oder auf den Rücken, und Lasten tragen wir auf der Schulter.

%%K t02-13-1 p
13. --- Im Bauch liegen der \VUgras{Magen} und der \VUgras{Darm}. Sie verdauen
unsere Speisen.

%%K t02-tit-3 sub
\VUsaut{4.4mm}
\VUpk{10.2pt}{40}{Die Gliedmaßen.}
\VUsaut{3.0mm}

%%K t02-14-1 p
14. --- Wir haben vier \VUgras{Gliedmaßen}: zwei \VUgras{Arme} und zwei
\VUgras{Beine}. Mit den Armen nehmen, halten, heben und sammeln wir Dinge auf; mit
den Beinen stehen, gehen und laufen wir.

%%K t02-15-1 p
15. --- Die \VUgras{Schulter} verbindet den Arm mit dem Rumpf. Ein sehr kräftiger
\VUgras{Muskel}, der Bizeps, bewegt den Arm, den der \VUgras{Ellbogen} mit dem
\VUgras{Unterarm} verbindet. Das \VUgras{Handgelenk} bildet das Gelenk der
\VUgras{Hand}, die in den \VUgras{Fingern} und den \VUgras{Nägeln} endet. Auf der
Hand erkennen wir eine \VUgras{Ader} (und eine Schlagader), in der das
\VUgras{Blut} fließt.

%%K t02-16-1 p
16. --- Die Teile des Beins sind: der \VUgras{Oberschenkel}, das \VUgras{Knie} und
die \VUgras{Kniekehle}, die \VUgras{Wade}, deren \VUgras{Fleisch} von der
\VUgras{Haut} bedeckt ist, der \VUgras{Knöchel} und der \VUgras{Fuß}. Die
\VUgras{Knochen} des Beins sind sehr dick und sehr fest. Am Ende des Fußes sitzen
die \VUgras{Zehen}. Die übrigen Teile sind die \VUgras{Sohle} und die
\VUgras{Ferse}.

%%K t02-17-1 p
17. --- Das ist eine fast vollständige Beschreibung des menschlichen Körpers.

%%K t02-17-2 p
\textit{Anmerkung.} --- \textit{Mit dem Satz „mir tut der ... weh (der Kopf, und
so weiter)“ kann der Lehrer diesen ganzen Wortschatz noch einmal durchgehen,
diesmal unter dem Gesichtspunkt der guten oder schlechten} \VUgras{Gesundheit}.

%%K t02-tit-4 sub
\VUsaut{5.0mm}
\VUpk{10.2pt}{40}{Turnen.}
\VUsaut{2.4mm}
\VUcentre{10.2pt}{0}{\textit{(Erzählt von einem der Schüler.)}}
\VUsaut{3.0mm}

%%K t02-18-1 p
18. --- Damit wir gelenkig, flink und gesund bleiben, müssen wir Beine und Arme
üben. Deshalb spielen wir und \VUgras{turnen} wir.

%%K t02-19-1 p
19. --- Unter der Woche finden diese beiden Übungen auf dem Schulhof statt (siehe
Tafel Nr. 1). Donnerstags und sonntags aber gehen wir hinaus auf eine große Wiese,
wo wir Platz zum Laufen und zum Vergnügen haben.

%%K t02-20-1 p
20. --- Unsere Lehrer und unsere Eltern kommen manchmal mit; die Übungen leitet
aber der \VUgras{Turnlehrer} \textsuperscript{(12)}.

%%K t02-21-1 p
21. --- Auf der Wiese, links vom Eingang, stehen die \VUgras{Turngeräte}
\textsuperscript{(1)}: wir klettern die \VUgras{Leiter} \textsuperscript{(2)} hinauf,
wir schwingen kopfüber an den \VUgras{Ringen} \textsuperscript{(3)} und am
\VUgras{Trapez} \textsuperscript{(4)}, wir ziehen uns Hand über Hand bis an das
obere Ende der \VUgras{Strickleiter} \textsuperscript{(5)} oder des
\VUgras{Knotentaus} \textsuperscript{(6)}.

%%K t02-22-1 p
22. --- Unterdessen springen unsere Freunde über das \VUgras{Sprungpferd}
\textsuperscript{(7)} oder über den \VUgras{Barren} \textsuperscript{(11)}. Andere
\VUgras{boxen} \textsuperscript{(8)} oder \VUgras{fechten} \textsuperscript{(9)} mit
Maske und \VUgras{Florett}. Wieder andere stemmen \VUgras{Hanteln}
\textsuperscript{(10)}.

%%K t02-tit-5 sub
\VUsaut{4.6mm}
\VUpk{10.2pt}{40}{Spiele.}
\VUsaut{3.0mm}

%%K t02-23-1 p
23. --- Rund um die Turner spielen Jungen und Mädchen die verschiedensten
\VUgras{Spiele}. Die Mädchen spielen mit ihrem \VUgras{Reifen}
\textsuperscript{(14)}; sie \VUgras{springen Seil} \textsuperscript{(13)} oder laden
ihre Brüder und Vettern zu einer Partie \VUgras{Krocket} \textsuperscript{(15)} ein.
Es macht ihnen große Freude, die \VUgras{Kugel} des Gegners mit dem
\VUgras{Holzschlägel} wegzuschlagen, gerade wenn er glaubt, er komme mühelos durch
das Tor!

%%K t02-24-1 p
24. --- Die Jungen mögen kräftigere Spiele lieber. Gern \VUgras{kegeln}
\textsuperscript{(16)} sie und werfen die Kegel mit der \VUgras{Kugel}
\textsuperscript{(17)} sauber um. Sie verschmähen auch den \VUgras{Kreisel}
\textsuperscript{(18)} und den \VUgras{Fangbecher} \textsuperscript{(19)} nicht, ja
nicht einmal das \VUgras{Tonnenspiel} \textsuperscript{(20)}. Am liebsten aber
spielen sie \VUgras{Murmeln} \textsuperscript{(21)} und \VUgras{Wandball}
\textsuperscript{(22)}, weil die Mauer des \VUgras{Gewächshauses}
\textsuperscript{(26)} sich zum Abprallen des leichten Balls hervorragend eignet.

%%K t02-25-1 p
25. --- Der kleine Johann auf seinen \VUgras{Stelzen} \textsuperscript{(24)} ist sehr
geschickt und hält das Gleichgewicht sehr gut. In seiner Nähe spielen ein paar
Freunde in diesem Gewächshaus und hinter der \VUgras{Hecke}
\VUgras{Verstecken} \textsuperscript{(25)}. Andere spielen lieber
\VUgras{Handklatschen} \textsuperscript{(28)} oder \VUgras{Federball}
\textsuperscript{(29)}, der von einem \VUgras{Schläger} zum anderen fliegt.

%%K t02-noto-1 noto
\VUnotes{143.2mm}{%
(*) Kinderspiele tragen oft rein nationale Namen. Wir haben sie im Ido so gut wie
möglich benannt, entweder nach der Handlung, die sie kennzeichnet, oder indem wir
den Namen aus dem einen oder anderen Land übernommen haben, wo er nicht zu
irreführend ist. Zum Beispiel: das \VUgras{Tonnenspiel} (deutsch und
französisch) heißt englisch \VUgras{toad in the hole} \textsuperscript{(20)}; die
Spieler versuchen dabei, ihre runden Scheiben durch das Maul der metallenen
Kröte zu werfen, die mit offenem Rachen auf der durchlöcherten Tonne sitzt. Der
\VUgras{Schulterbock} \textsuperscript{(30)} unterscheidet sich vom
\VUgras{Bockspringen} nur dadurch, dass die Springer sich beim Sprung über den
Kopf auf die Schultern des Partners stützen, während sie beim
„\VUgras{Bockspringen}“ nur den Rücken berühren. --- Oder auch: ein Teil der
Spieler bückt sich, während die anderen ihnen mit den Händen auf den Schultern
über den Rücken springen; die einen müssen den Stoß aushalten, ohne
nachzugeben, die anderen springen, ohne das Gleichgewicht zu verlieren (siehe die
Tafel selbst).%
}

%%K t02-26-1 p
26. --- In der Mitte der Wiese stehen zwei Bäume. Am Fuß des einen haben sieben
oder acht Jungen eine Runde \VUgras{Schulterbock} \textsuperscript{(30)} (*)
angefangen. Dieses Spiel kennt man nicht in jedem Land; was aber
\VUgras{Bockspringen} ist, weiß jeder, und gerade jetzt spielen die Kinder es
nicht.

%%K t02-tit-6 sub
\VUsaut{5.0mm}
\VUpk{10.2pt}{40}{Spiele im Freien.}
\VUsaut{3.4mm}

%%K t02-27-1 p
27. --- Auch \VUgras{Blindekuh} \textsuperscript{(31)} macht großen Spaß; Mädchen und
Jungen binden sich abwechselnd die \VUgras{Augenbinde} um und fangen die
Ungeschickten, die sich fangen lassen.

%%K t02-28-1 p
28. --- Mitten auf der Wiese wird ein sehr französisches Spiel gespielt, ein etwas
grobes allerdings: der \VUgras{Bär} \textsuperscript{(32)}, viel lauter als das
friedliche Spiel mit dem \VUgras{Drachen} \textsuperscript{(33)} oder gar als das
englische \VUgras{Tennis} \textsuperscript{(34)}.

%%K t02-29-1 p
29. --- Dieser letzte Sport ist die Freude der älteren Schüler. Sogar unsere Lehrer
machen gern mit. Er verlangt viel Geschick und Beweglichkeit, und ich liebe ihn.
Die \VUgras{Schaukel} \textsuperscript{(35)} überlasse ich den Mädchen.

%%K t02-30-1 p
30. --- Auf ein anderes englisches Spiel, das leicht gefährlich werden kann, bin
ich auch nicht scharf.

%%K t02-30-1b p
Ich meine \VUgras{Fußball} \textsuperscript{(36)}, die Freude meines älteren
Bruders. Mir ist unser altes \VUgras{Barlauf} \textsuperscript{(38)} lieber, genauso
gesund und viel weniger brutal.

%%K t02-tit-7 sub
\VUsaut{4.6mm}
\VUpk{10.2pt}{40}{Radfahren und Ballspiele.}
\VUsaut{3.0mm}

%%K t02-31-1 p
31. --- Ich fahre auch gern mit dem \VUgras{Fahrrad} \textsuperscript{(37)} um die
Wiese herum.

%%K t02-32-1 p
32. --- Nächstes Jahr lerne ich \VUgras{reiten} \textsuperscript{(39)}. Ein Pferd ist
ein so schöner Anblick, wenn es anmutig schreitet oder trabt und im Galopp über
die Hindernisse setzt!

%%K t02-33-1 p
33. --- Im Sommer lernen wir \VUgras{schwimmen} \textsuperscript{(40)}. Wir tauchen
und baden mit Vergnügen im klaren, kühlen Wasser des Flusses. Manchmal lassen wir
zum Schluss einen \VUgras{Papierballon} \textsuperscript{(41)} steigen, der
majestätisch in die Luft aufsteigt, sobald er mit heißer Luft gefüllt ist.

%%K t02-34-1 p
34. --- Es gibt noch viele andere Spiele, aber ich käme mit dem Aufzählen nie zu
Ende, und man sieht schon an diesem kurzen Bericht, dass die Donnerstage auf
unserer Wiese alles andere als langweilig sind.

%%K t02-34-2 p
N.B. --- \textit{Der Lehrer kann dieses Kapitel leicht erweitern, indem er die
wichtigeren Spiele einzeln genauer beschreibt.}
\VUsaut{6.0mm}
\VUfilet{22mm}
